\documentclass[11pt]{article}

\usepackage[margin=1in]{geometry}
\usepackage{amsmath}
\usepackage{amssymb}

\usepackage{graphicx}
\usepackage{bm}
\usepackage{booktabs}
\usepackage{array}
\usepackage{amsthm}
\usepackage{enumerate}
\usepackage[numbers,sort&compress]{natbib}
\usepackage{hyperref}
\usepackage{lineno}

\hypersetup{
  colorlinks=true,
  linkcolor=blue,
  citecolor=blue,
  urlcolor=blue
}

\newtheorem{theorem}{Theorem}
\newtheorem{proposition}{Proposition}
\newtheorem{corollary}{Corollary}
\newtheorem{lemma}{Lemma}

\begin{document}

\title{A Notational Reformulation of Torus-Knot Self-Consistency in Flow--Constraint Variables}

\author{Steve Bico Mujjabi, MD\thanks{msbico@gmail.com}\\
VuraLabs, Kampala, Uganda\\
Independent researcher\\
ORCID: \href{https://orcid.org/0009-0001-0556-5516}{0009-0001-0556-5516}\\
\url{https://github.com/bampita-bico/CDFD-Part-I-Release}}

\date{July 2026}

\maketitle

\begin{abstract}
This paper records a flow--constraint notation for the torus-knot
self-consistency construction used in the earlier CDFD papers. It does not
derive a physical vacuum equation of state or establish the physical origin of
the twist condition. The relation $\Psi_s=\Phi/C$ is postulated notation whose
utility must be judged by an independent physical prediction.

The paper maps selected Fourier quantities onto three CDFD labels: Flow,
Constraint, and Adaptive Operating Ratio. The resulting equality is an
algebraic restatement of those choices, not an independent result. The
paper-local supplementary scripts reproduce the algebra and figures; they do
not validate the physical model.
\end{abstract}

\noindent\textbf{Keywords:} Constraint-Driven Field Theory, Constraint-Driven Flux Dynamics, adaptive operating ratio, vacuum structure, Mujjabi tests, reproducible physics, vacuum balance, public CDFD model

\vspace{0.75em}
\linenumbers

\section*{Project Context}
This manuscript constitutes Part I (Paper 10) of the multi-part series
\textit{Constraint-Driven Flux Dynamics (CDFD)}. The underlying mathematical
architecture builds directly upon the concepts established in Paper 09 of this
series~\cite{MujjabiIX2026}.

\section*{Symbol Conventions}

\noindent\textbf{CDFD physical-vacuum symbol convention.}
\begin{center}
{\small
\begin{tabular}{p{0.16\linewidth}p{0.33\linewidth}p{0.41\linewidth}}
\hline
Symbol & Part I meaning & Physical reading \\
\hline
$J$ or $\Phi$ & Driving flux or throughput & Transport current, field throughput, circulation drive, or generic flux; a subscript fixes the exact channel \\
$C$ & Active constraint or capacity burden & Vacuum transport capacity, geometric confinement, stiffness, boundary resistance, or load-bearing limit \\
$S$ & Surface responsiveness & Local ability of the vacuum or modeled surface to reconfigure under drive \\
$M_s$ & Structural memory & Retained geometry, phase history, stiffness history, or accumulated constraint state \\
$\Psi_s$ & Adaptive operating ratio $(J/C)\cdot S\cdot M_s$ or $(\Phi/C)\cdot S\cdot M_s$ & Local bookkeeping gauge for constrained operation; it is not a new universal constant \\
$R,a,\chi$ & Ring radius, core radius, and aspect ratio $R/a$ & Geometry used to connect vortex structure to dimensionless electromagnetic scales \\
$\rho_0,c,\kappa$ & Vacuum density scale, propagation speed, and transport stiffness & Medium parameters used in stability, pressure, and self-consistency calculations \\
\hline
\end{tabular}
}
\end{center}

\noindent\textbf{Greek-symbol convention.}
\begin{center}
{\small
\begin{tabular}{p{0.16\linewidth}p{0.33\linewidth}p{0.41\linewidth}}
\hline
Symbol & Local role & Interpretation rule \\
\hline
$\alpha$ & Fine-structure constant unless a local equation states otherwise & In the Part I physics papers, un-subscripted $\alpha$ means the electromagnetic fine-structure constant \\
$\beta,\gamma$ & Local response, stiffness, coupling, or correction coefficients & Equation-local coefficients; subscripts bind them to the named mode, channel, or approximation \\
$\theta,\phi,\Phi_c$ & Phase, orientation, or chirality angles & Used for torus-knot orientation, vacuum phase, or chiral phase geometry as defined in the nearby equation \\
$\delta,\epsilon$ & Perturbation, residual, tolerance, or small-offset scales & Local stability margins, numerical residuals, or experimental tolerances \\
$\lambda,\mu,\nu$ & Rate, mass, mode, or frequency labels & Meaning is fixed by the equation, subscript, and physics context where the symbol appears \\
$\rho,\sigma,\tau$ & Density, spread/noise, or time-scale terms & Local density, variance, stress, relaxation, or transport-memory quantities \\
$\Omega,\Lambda$ & State/domain space or integrated measure where defined & Reserved for explicitly defined state spaces, spectra, or synthesis-level quantities \\
\hline
\end{tabular}
}
\end{center}


\section{Introduction}

Across Papers~I--IX~\cite{MujjabiI2026,MujjabiII2026,MujjabiIII2026,%
MujjabiIV2026,MujjabiV2026,MujjabiVI2026,MujjabiVII2026,MujjabiVIII2026,MujjabiIX2026},
CDFD introduced a torus-knot parametrisation and several internal algebraic
identities. \texttt{CORRECTION\_STATUS\_2026-08-18.md} records the binding distinction
between those conditional identities and the unsupported physical mapping.

The entire predictive power of the higher families rests on the unique solution
to the self-consistency condition:
\begin{equation}
  n\theta_n \;=\; \widetilde{Q}_n(\theta_n) \;\equiv\; \frac{\sum_{k=0}^{n-1} A_k^2}{\bigl(\sum_{k=0}^{n-1} |A_k|\bigr)^2}\,.
  \label{eq:self_consistency}
\end{equation}
In Paper~V, for $n=3$, this was introduced as a postulate matching the
chirality breaking angle to the algebraic constant $Q_3 = 2/3$.
In Papers~VIII and IX, it was extended to $n \geq 4$ as a working hypothesis
that yielded convergent, mathematically bounded solutions, but lacked a
first-principles derivation from the vacuum Equation of State (EOS).

This paper only rewrites the torus-knot construction in the public CDFD
flow--constraint notation; it does not close the open physical problem.


\section{The Tripartite Vacuum Framework}

The CDFD vacuum is modeled as an elastic, transport-capacity-limited continuous medium,
in the spirit of Dirac's physical vacuum~\cite{Dirac1951} and
Volovik's condensate analogy~\cite{Volovik2003}.
The vortex dynamics rest on Helmholtz's conservation
theorems~\cite{Helmholtz1858} and the Faddeev--Niemi knot
soliton framework~\cite{FaddeevNiemi1997}.
For this manuscript, the following postulated notational identity is adopted:
\begin{equation}
  \Psi_s \;=\; \frac{\Phi}{C}
  \label{eq:UE_axiom}
\end{equation}
where $\Phi$ is flow or activity, $C$ is constraint, resistance, or capacity,
and $\Psi_s$ is the resulting adaptive operating ratio.

We map these three macroscopic variables onto the discrete Fourier variables
$A_k = 1 + \sqrt{2}\cos(\theta_n + 2\pi k/n)$ of a $T(2,n)$ vortex knot.

\subsection{Flow ($\Phi$): The Field Energy}
The total energy or ``activity'' of the field configuration is proportional
to the $L^2$ norm (the sum of the squares) of the amplitudes. By the Fourier
identity of Paper~VI, this is invariant under orientation $\theta_n$:
\begin{equation}
  \Phi \;\propto\; \sum_{k=0}^{n-1} A_k^2 \;=\; 2n\,.
  \label{eq:flow_phi}
\end{equation}

\begin{figure}[ht]
  \centering
  \includegraphics[width=0.8\textwidth]{figures/fig10_0_panel.pdf}
  \caption{The flow energy ratio $\Phi/n$ is exactly $2.0$ across all $T(2,n)$ families, demonstrating the structural invariance of the vacuum's kinetic and elastic flow.}
  \label{fig:phi_ratio}
\end{figure}

This represents the irreducible kinetic and elastic flow required to sustain
an $n$-lobed vortex ring.

\subsection{Constraint ($C$): The Transport-Capacity Limit}
A central postulate of CDFD (Paper~I) is that the vacuum possesses a finite
transport capacity, which dictates the electromagnetic fine-structure constant
$\alpha \equiv e^2/(\hbar c)$.
For a multi-lobed knot, the structural stress on the medium is dictated by the
total absolute flux. In the discrete basis, the total flux is the $L^1$ norm
$\sum |A_k|$.
Because the constraint scales with the cross-sectional area of the flux tubes
interacting in the medium, the total resistance $C$ scales as the square of the flux:
\begin{equation}
  C \;\propto\; \left(\sum_{k=0}^{n-1} |A_k|\right)^2\,.
  \label{eq:constraint_C}
\end{equation}
This term varies strongly with $\theta_n$, as the lobes interfere constructively
and destructively with the vacuum background.

\subsection{Equilibrium ($\Psi_s$): The Geometric Twist}
When the knot is embedded in the vacuum, the mismatch between the continuous
chiral background and the discrete $n$-fold symmetry of the knot induces a
geometric twist.
Let $\theta_n$ be the phase twist per lobe. The total equilibrium deformation
(or total accumulated twist) of the configuration is:
\begin{equation}
  \Psi_s \;=\; n\theta_n\,.
  \label{eq:equilibrium_psi}
\end{equation}


\section{The $\Psi_s=\Phi/C$ Notational Identity}

Substituting the chosen CDFD mappings \eqref{eq:flow_phi},
\eqref{eq:constraint_C}, and \eqref{eq:equilibrium_psi} into the postulated
notation $\Psi_s = \Phi / C$ gives:
  \begin{equation}
    n\theta_n \;=\; \frac{\sum_{k=0}^{n-1} A_k^2}{\left(\sum_{k=0}^{n-1} |A_k|\right)^2}\,.
  \end{equation}
The right-hand side is the chosen definition of the modified Koide ratio
$\widetilde{Q}_n(\theta_n)$ from Eq.~\eqref{eq:self_consistency}. This is a
consistent relabelling of the construction, not a theorem, proof, or physical
derivation.

\begin{figure}[ht]
  \centering
  \includegraphics[width=0.8\textwidth]{figures/fig10_1_panel.pdf}
  \caption{The chosen labels $\Psi_s$ and $\Phi/C$ plotted across the stated torus-knot construction. The diagonal identity verifies the internal algebraic relabelling only.}
  \label{fig:psi_vs_phi_over_c}
\end{figure}

\begin{figure}[ht]
  \centering
  \includegraphics[width=0.8\textwidth]{figures/fig10_2_panel.pdf}
  \caption{Absolute residuals $|\Phi / C - \Psi_s|$. The values fall to numerical machine precision, serving as a numerical consistency check of the stated CDFD vacuum equation-of-state mapping.}
  \label{fig:residuals}
\end{figure}

\subsection{Interpretation boundary}
The mapping does not show that the observed Koide ratio is a vacuum equation of
state, nor that the higher-knot quantities are particle spectra. Any physical
interpretation requires an independent mapping and a test that outperforms
standard theory.

\section{Vacuum Engineering as a Hypothetical Control Programme}

The public balance law is also the first engineering form of the theory. Vacuum
engineering, in this manuscript, does not mean a claimed device. It means a
control program: identify which variable in the balance law can be perturbed,
then measure whether the response follows the CDFD operator better than the
standard model for that domain.

This is why vacuum engineering is introduced here as a future field of study:
the disciplined attempt to manipulate, bound, or falsify the state variables
$\Phi$, $C$, $S$, and $M_s$ in systems where the vacuum or a vacuum-analogue
medium is being driven near a transport threshold. The field begins with
protocols and null tests, not devices. Its first deliverables are calibrated
phase-residual limits, spectral-residual limits, hysteresis bounds, and
constraint-threshold maps.

\begin{center}
\small
\begin{tabular}{p{0.16\textwidth}p{0.29\textwidth}p{0.28\textwidth}p{0.15\textwidth}}
\toprule
\textbf{Variable} & \textbf{Engineering handle} & \textbf{Observable} & \textbf{Status}\\
\midrule
$\Phi$ or $J$ & field intensity, current, shear, angular momentum, density gradient & phase, spectral, transport, or force response & controllable in testbeds\\
$C$ & geometry, boundary, topology, density, confinement & capacity threshold or saturation point & model-dependent\\
$S$ & susceptibility or response surface & slope, compliance, or relaxation curve & open variable\\
$M_s$ & prior loading or pulse history & hysteresis, lag, or recovery time $\tau_M$ & falsification target\\
\bottomrule
\end{tabular}
\end{center}

We call this the Mujjabi Vacuum Engineering Principle: steer $\Phi$, $C$, $S$,
or $M_s$ to move a system toward, across, or away from $\Psi_s=1$, then measure
the residual. Paper~XII converts this principle into named laboratory tests.

\subsection{CDFD Runtime experiment hooks}

The root CDFD Runtime now contains experiment scripts that turn the above
principle into machine-readable candidate outputs. These are not empirical
confirmations; they are public stress tests that help decide which claims are
worth carrying into laboratory protocols.

\begin{center}
\small
\begin{tabular}{p{0.28\textwidth}p{0.40\textwidth}p{0.22\textwidth}}
\toprule
\textbf{Runtime output} & \textbf{Observed model-side result} & \textbf{Release use}\\
\midrule
\texttt{discovery\_vacuum\_hysteresis.h5} &
central $M_s$ peaked at $15.58$ after high-flux pulsing &
candidate vacuum-memory protocol\\
\texttt{discovery\_knot\_n7.h5} &
$T(2,7)$ stability index drifted from $0.992686$ to $0.987976$ &
higher-knot stress test\\
\texttt{discovery\_ool\_phase\_results.h5} &
a $50\times50$ sweep found four grid points near $\Lambda=1$ &
cross-Part-II bridge only\\
\texttt{discovery\_surface\_evolution.h5} &
$\Psi_s$ ran from $5.15$ to $978.99$ under overload &
instability boundary, not support\\
\texttt{frontier\_sweep.h5} &
CODATA 2022 $\chi^*=137.035999177$ used across 2500 points &
hypothesis triage\\
\bottomrule
\end{tabular}
\end{center}

The publication rule is strict: a runtime file may motivate a paper statement
only when the manuscript names the state variables, records the parameters, and
states the falsification condition. The refreshed reports in
\texttt{experiments/reports/} are therefore simulation logs, not evidence that
any vacuum-engineering device or clinical intervention already exists.


\section{Scope Boundary and Summary}

This paper makes no cross-scale biology, medicine, ecology, or Fibonacci claim.
Such analogies were removed because formal resemblance does not supply a
physical derivation or incremental prediction. The remaining content is a
notational restatement of a selected torus-knot construction. It does not close
the Part I physical programme, establish a mass hierarchy, or support
downstream Parts II and III as a source of established infrastructure.


\section*{Supplementary Material}

This manuscript is accompanied by release-archive supplementary material in
the canonical Part I tree. The relevant paper-local Python scripts, numerical checks, generated figure outputs,
tabulated data, figure assets, and environment notes are maintained under
\texttt{notebooks/}, \texttt{outputs/}, \texttt{figures/}, and
\texttt{REPRODUCIBILITY.md} in \texttt{Part\_I\_Fundamental\_Physics/}.


\section*{Data and Code Availability}

The release-archive supplementary material is included in the archived Part I
release on Zenodo at \url{https://doi.org/10.5281/zenodo.20250820}. The current version DOI is assigned by Zenodo after each GitHub release. Zenodo is the
permanent release archive, and the public source archive is available on GitHub at
\url{https://github.com/bampita-bico/CDFD-Part-I-Release}. The broader public CDFD
Runtime is separate reusable software infrastructure; the numerical values reported
here are reproduced from this paper's Supplementary Material.


\section*{Acknowledgments}
The author used AI-assisted tools for formatting checks and code-verification
support. The scientific claims, interpretations, and final manuscript decisions
are the author's responsibility.

The author is indebted to the foundational work of
Lord Kelvin~\cite{Thomson1867} (vortex atom hypothesis),
P.A.M.~Dirac~\cite{Dirac1951} (vacuum as physical medium),
G.E.~Volovik~\cite{Volovik2003} (vacuum as condensate),
H.~von Helmholtz~\cite{Helmholtz1858} (vortex theorems),
L.~Faddeev and A.J.~Niemi~\cite{FaddeevNiemi1997} (knots and particles),
Y.~Koide~\cite{Koide1983} (lepton mass formula),
and C.A.~Brannen~\cite{Brannen2006} (Koide parametrisation)
--- whose ideas this series has attempted to unify within a single
geometric framework rooted in the axiom $\Psi_s = \Phi / C$.

\section*{Declarations}

\subsection*{Funding}
This research received no external institutional funding.

\subsection*{Author Contributions}
Steve Bico Mujjabi: conceptualization, methodology, formal analysis, software,
validation, visualization, writing--original draft, and writing--review and
editing.

\subsection*{Conflicts of Interest}
The author declares no conflicts of interest.

\section*{References}
\bibliographystyle{plainnat}
\bibliography{references}

\end{document}
